TEV assigns 153 occupations to tier 1, 284 to tier 2, and 65 to tier 3. Table~\ref{tab:v1_jev_job_tiers} compares tier 1 with pooled tiers 2 and 3. In the preferred within-CNO1 specification, the unemployment coefficients are 0.065 log points during the adjustment period and 0.116 later; both are statistically significant at the 1 percent level. The contract coefficients are $-0.086$ and $-0.083$ log points, with only the adjustment-period estimate significant at the 5 percent level. Excluding 2021 yields unemployment coefficients of 0.042 and 0.093, both significant at the 5 percent level, and contract coefficients of $-0.057$ and $-0.055$, neither significant at that level.

The preferred pre-treatment joint-null tests reject for both unemployment ($p=0.010$) and contracts ($p=0.007$). Without 2021, the unemployment test continues to reject ($p=0.005$), while the contract test does not ($p=0.226$). The tier assignments and the resulting diagnostics therefore matter for interpretation: these estimates describe differential occupational adjustment and do not identify a protective causal effect of bundling. Figures~\ref{fig:v1_jev_tiers_unemployed} and~\ref{fig:v1_jev_tiers_contracts} show the corresponding event-study paths.

Table~\ref{tab:v1_job_tier_mediation} adds tier-by-year-month fixed effects to the baseline exposure specification, allowing each of the three TEV tiers its own monthly path. The unemployment coefficients decline from 0.012 and 0.017 to 0.008 and 0.010; neither adjusted estimate is significant at the 5 percent level. For contracts, they change from 0.014 and 0.013 to 0.017 and 0.014, with only the adjustment-period coefficient significant at the 5 percent level. Conditioning on tier dynamics changes the exposure gradient, but does not identify which tier or economic mechanism accounts for the change.
